% Seminar 4 - Pagina 3. Fragment pentru
% % Seminar 4 - Pagina 3. Fragment pentru
% % Seminar 4 - Pagina 3. Fragment pentru
% % Seminar 4 - Pagina 3. Fragment pentru
% \input{seminars/pages/seminar4-pagina3-ro}; fără preambul sau pagină nouă.
% NOTE EDITORIALE:
% - La exercițiul 7, manuscrisul pare să scrie "nu e inj.".
%   Incluziunea canonică Z -> Q ESTE injectivă, însă nu este surjectivă.
%   Păstrăm lectura aparentă, marcată "ca în manuscris", fără corectare
%   tacită.
% - Monoizii din exercițiul 7 sunt priviți cu înmulțirea.
%   Diagrama arată definiția epimorfismului prin egalitatea a două
%   morfisme din Q care coincid pe Z.
% - În calculul pentru v, manuscrisul folosește o linie lungă în locul
%   repetării pașilor; păstrăm această prescurtare.
% - În exercițiul 8 se presupune un izomorfism ipotetic: funcția este
%   bijectivă și respectă produsul, iar la final se folosește injectivitatea.
% - Notația de lângă acolada pentru f(a)f(b)=4 are un al doilea rând
%   greu de citit (f(b)=...?); f(a)=1 => a=1 este clar.
% - "Numere prime" se referă la numere prime în N^*.
% - Textul șters și corecturile barate nu sunt reconstruite.

\subsection*{Seminar 4 - Pagina 3}

\noindent\underline{Exercițiul 7.}
Incluziunea
\[
 i:(\mathbb Z,\cdot)\longrightarrow(\mathbb Q,\cdot),
 \qquad i(x)=x,
\]
este epimorfism de monoizi care \emph{nu e injectiv}
(așa cum apare în manuscris).

\smallskip
\noindent\underline{Epimorfism.}
Pentru orice monoid $M$ și orice morfisme de monoizi $u,v$
astfel încât $u\circ i=v\circ i$, avem $u=v$.
\[
\begin{tikzcd}[column sep=large]
 \mathbb Z \arrow[r,"i"] &
 \mathbb Q \arrow[r,shift left=.7ex,"u"]
            \arrow[r,shift right=.7ex,swap,"v"] & M
\end{tikzcd}
\]

\noindent Fie $(M,\cdot)$ monoid și $u,v:\mathbb Q\to M$
morfisme de monoizi cu $u\circ i=v\circ i$. Atunci
\[
 u(m)=v(m)\qquad(\forall m\in\mathbb Z).
\]
Vrem să arătăm că $u(r)=v(r)$ pentru orice $r\in\mathbb Q$.
Fie $m\in\mathbb Z$, $m\ne0$. Avem
\[
\begin{aligned}
 e=u(1)
  &=u\!\left(\frac{m}{m}\right)
   =u\!\left(m\cdot\frac1m\right)
   =u(m)u\!\left(\frac1m\right)
   =u\!\left(\frac1m\right)u(m),\\
 e=v(1)
  &=\cdots
   =v\!\left(\frac1m\right)v(m).
\end{aligned}
\]
Deci $u(1/m)$ este inversul lui $u(m)$, iar $v(1/m)$ este
inversul lui $v(m)$. Cum $u(m)=v(m)$, rezultă
\[
 u\!\left(\frac1m\right)=v\!\left(\frac1m\right).
\]
Pentru $r=n/m\in\mathbb Q$, cu $m\ne0$,
\[
 u(r)=u(n)u\!\left(\frac1m\right)
     =v(n)v\!\left(\frac1m\right)=v(r).
\]
Prin urmare $u=v$.

\smallskip
\noindent [Orice epimorfism de grupuri este surjectiv.]

\medskip
\noindent\underline{Exercițiul 8.}
Monoizii
\[
 (\mathbb N^*,\cdot),\qquad
 (3\mathbb N+1,\cdot),\qquad
 3\mathbb N+1=\{3k+1\mid k\in\mathbb N\},
\]
nu sunt izomorfi.

\noindent\underline{Demonstrație.}
Presupunem că există un izomorfism
\[
 f:\mathbb N^*\longrightarrow 3\mathbb N+1,
 \qquad f(xy)=f(x)f(y)
 \quad(\forall x,y\in\mathbb N^*).
\]
Din surjectivitatea lui $f$ rezultă că există $x_1,x_2,x_3\in\mathbb N^*$
astfel încât
\[
 f(x_1)=4,\qquad f(x_2)=10,\qquad f(x_3)=25.
\]
Arătăm că $x_1,x_2,x_3$ sunt numere prime.
Dacă $x_1=ab$, cu $a,b>1$, atunci
\[
 f(x_1)=f(a)f(b)=4
 \quad\Longrightarrow\quad f(a)=1
 \quad\Longrightarrow\quad a=1,
 \qquad\text{fals.}
\]
Analog, $x_2$ și $x_3$ sunt numere prime.
Dar $4\cdot25=10^2$, deci
\[
\begin{aligned}
 f(x_1)f(x_3)&=f(x_2)^2,\\
 f(x_1x_3)&=f(x_2^2),\\
 x_1x_3&=x_2^2
  \quad\Longrightarrow\quad x_1=x_2=x_3,
  \qquad\text{fals.}
\end{aligned}
\]
; fără preambul sau pagină nouă.
% NOTE EDITORIALE:
% - La exercițiul 7, manuscrisul pare să scrie "nu e inj.".
%   Incluziunea canonică Z -> Q ESTE injectivă, însă nu este surjectivă.
%   Păstrăm lectura aparentă, marcată "ca în manuscris", fără corectare
%   tacită.
% - Monoizii din exercițiul 7 sunt priviți cu înmulțirea.
%   Diagrama arată definiția epimorfismului prin egalitatea a două
%   morfisme din Q care coincid pe Z.
% - În calculul pentru v, manuscrisul folosește o linie lungă în locul
%   repetării pașilor; păstrăm această prescurtare.
% - În exercițiul 8 se presupune un izomorfism ipotetic: funcția este
%   bijectivă și respectă produsul, iar la final se folosește injectivitatea.
% - Notația de lângă acolada pentru f(a)f(b)=4 are un al doilea rând
%   greu de citit (f(b)=...?); f(a)=1 => a=1 este clar.
% - "Numere prime" se referă la numere prime în N^*.
% - Textul șters și corecturile barate nu sunt reconstruite.

\subsection*{Seminar 4 - Pagina 3}

\noindent\underline{Exercițiul 7.}
Incluziunea
\[
 i:(\mathbb Z,\cdot)\longrightarrow(\mathbb Q,\cdot),
 \qquad i(x)=x,
\]
este epimorfism de monoizi care \emph{nu e injectiv}
(așa cum apare în manuscris).

\smallskip
\noindent\underline{Epimorfism.}
Pentru orice monoid $M$ și orice morfisme de monoizi $u,v$
astfel încât $u\circ i=v\circ i$, avem $u=v$.
\[
\begin{tikzcd}[column sep=large]
 \mathbb Z \arrow[r,"i"] &
 \mathbb Q \arrow[r,shift left=.7ex,"u"]
            \arrow[r,shift right=.7ex,swap,"v"] & M
\end{tikzcd}
\]

\noindent Fie $(M,\cdot)$ monoid și $u,v:\mathbb Q\to M$
morfisme de monoizi cu $u\circ i=v\circ i$. Atunci
\[
 u(m)=v(m)\qquad(\forall m\in\mathbb Z).
\]
Vrem să arătăm că $u(r)=v(r)$ pentru orice $r\in\mathbb Q$.
Fie $m\in\mathbb Z$, $m\ne0$. Avem
\[
\begin{aligned}
 e=u(1)
  &=u\!\left(\frac{m}{m}\right)
   =u\!\left(m\cdot\frac1m\right)
   =u(m)u\!\left(\frac1m\right)
   =u\!\left(\frac1m\right)u(m),\\
 e=v(1)
  &=\cdots
   =v\!\left(\frac1m\right)v(m).
\end{aligned}
\]
Deci $u(1/m)$ este inversul lui $u(m)$, iar $v(1/m)$ este
inversul lui $v(m)$. Cum $u(m)=v(m)$, rezultă
\[
 u\!\left(\frac1m\right)=v\!\left(\frac1m\right).
\]
Pentru $r=n/m\in\mathbb Q$, cu $m\ne0$,
\[
 u(r)=u(n)u\!\left(\frac1m\right)
     =v(n)v\!\left(\frac1m\right)=v(r).
\]
Prin urmare $u=v$.

\smallskip
\noindent [Orice epimorfism de grupuri este surjectiv.]

\medskip
\noindent\underline{Exercițiul 8.}
Monoizii
\[
 (\mathbb N^*,\cdot),\qquad
 (3\mathbb N+1,\cdot),\qquad
 3\mathbb N+1=\{3k+1\mid k\in\mathbb N\},
\]
nu sunt izomorfi.

\noindent\underline{Demonstrație.}
Presupunem că există un izomorfism
\[
 f:\mathbb N^*\longrightarrow 3\mathbb N+1,
 \qquad f(xy)=f(x)f(y)
 \quad(\forall x,y\in\mathbb N^*).
\]
Din surjectivitatea lui $f$ rezultă că există $x_1,x_2,x_3\in\mathbb N^*$
astfel încât
\[
 f(x_1)=4,\qquad f(x_2)=10,\qquad f(x_3)=25.
\]
Arătăm că $x_1,x_2,x_3$ sunt numere prime.
Dacă $x_1=ab$, cu $a,b>1$, atunci
\[
 f(x_1)=f(a)f(b)=4
 \quad\Longrightarrow\quad f(a)=1
 \quad\Longrightarrow\quad a=1,
 \qquad\text{fals.}
\]
Analog, $x_2$ și $x_3$ sunt numere prime.
Dar $4\cdot25=10^2$, deci
\[
\begin{aligned}
 f(x_1)f(x_3)&=f(x_2)^2,\\
 f(x_1x_3)&=f(x_2^2),\\
 x_1x_3&=x_2^2
  \quad\Longrightarrow\quad x_1=x_2=x_3,
  \qquad\text{fals.}
\end{aligned}
\]
; fără preambul sau pagină nouă.
% NOTE EDITORIALE:
% - La exercițiul 7, manuscrisul pare să scrie "nu e inj.".
%   Incluziunea canonică Z -> Q ESTE injectivă, însă nu este surjectivă.
%   Păstrăm lectura aparentă, marcată "ca în manuscris", fără corectare
%   tacită.
% - Monoizii din exercițiul 7 sunt priviți cu înmulțirea.
%   Diagrama arată definiția epimorfismului prin egalitatea a două
%   morfisme din Q care coincid pe Z.
% - În calculul pentru v, manuscrisul folosește o linie lungă în locul
%   repetării pașilor; păstrăm această prescurtare.
% - În exercițiul 8 se presupune un izomorfism ipotetic: funcția este
%   bijectivă și respectă produsul, iar la final se folosește injectivitatea.
% - Notația de lângă acolada pentru f(a)f(b)=4 are un al doilea rând
%   greu de citit (f(b)=...?); f(a)=1 => a=1 este clar.
% - "Numere prime" se referă la numere prime în N^*.
% - Textul șters și corecturile barate nu sunt reconstruite.

\subsection*{Seminar 4 - Pagina 3}

\noindent\underline{Exercițiul 7.}
Incluziunea
\[
 i:(\mathbb Z,\cdot)\longrightarrow(\mathbb Q,\cdot),
 \qquad i(x)=x,
\]
este epimorfism de monoizi care \emph{nu e injectiv}
(așa cum apare în manuscris).

\smallskip
\noindent\underline{Epimorfism.}
Pentru orice monoid $M$ și orice morfisme de monoizi $u,v$
astfel încât $u\circ i=v\circ i$, avem $u=v$.
\[
\begin{tikzcd}[column sep=large]
 \mathbb Z \arrow[r,"i"] &
 \mathbb Q \arrow[r,shift left=.7ex,"u"]
            \arrow[r,shift right=.7ex,swap,"v"] & M
\end{tikzcd}
\]

\noindent Fie $(M,\cdot)$ monoid și $u,v:\mathbb Q\to M$
morfisme de monoizi cu $u\circ i=v\circ i$. Atunci
\[
 u(m)=v(m)\qquad(\forall m\in\mathbb Z).
\]
Vrem să arătăm că $u(r)=v(r)$ pentru orice $r\in\mathbb Q$.
Fie $m\in\mathbb Z$, $m\ne0$. Avem
\[
\begin{aligned}
 e=u(1)
  &=u\!\left(\frac{m}{m}\right)
   =u\!\left(m\cdot\frac1m\right)
   =u(m)u\!\left(\frac1m\right)
   =u\!\left(\frac1m\right)u(m),\\
 e=v(1)
  &=\cdots
   =v\!\left(\frac1m\right)v(m).
\end{aligned}
\]
Deci $u(1/m)$ este inversul lui $u(m)$, iar $v(1/m)$ este
inversul lui $v(m)$. Cum $u(m)=v(m)$, rezultă
\[
 u\!\left(\frac1m\right)=v\!\left(\frac1m\right).
\]
Pentru $r=n/m\in\mathbb Q$, cu $m\ne0$,
\[
 u(r)=u(n)u\!\left(\frac1m\right)
     =v(n)v\!\left(\frac1m\right)=v(r).
\]
Prin urmare $u=v$.

\smallskip
\noindent [Orice epimorfism de grupuri este surjectiv.]

\medskip
\noindent\underline{Exercițiul 8.}
Monoizii
\[
 (\mathbb N^*,\cdot),\qquad
 (3\mathbb N+1,\cdot),\qquad
 3\mathbb N+1=\{3k+1\mid k\in\mathbb N\},
\]
nu sunt izomorfi.

\noindent\underline{Demonstrație.}
Presupunem că există un izomorfism
\[
 f:\mathbb N^*\longrightarrow 3\mathbb N+1,
 \qquad f(xy)=f(x)f(y)
 \quad(\forall x,y\in\mathbb N^*).
\]
Din surjectivitatea lui $f$ rezultă că există $x_1,x_2,x_3\in\mathbb N^*$
astfel încât
\[
 f(x_1)=4,\qquad f(x_2)=10,\qquad f(x_3)=25.
\]
Arătăm că $x_1,x_2,x_3$ sunt numere prime.
Dacă $x_1=ab$, cu $a,b>1$, atunci
\[
 f(x_1)=f(a)f(b)=4
 \quad\Longrightarrow\quad f(a)=1
 \quad\Longrightarrow\quad a=1,
 \qquad\text{fals.}
\]
Analog, $x_2$ și $x_3$ sunt numere prime.
Dar $4\cdot25=10^2$, deci
\[
\begin{aligned}
 f(x_1)f(x_3)&=f(x_2)^2,\\
 f(x_1x_3)&=f(x_2^2),\\
 x_1x_3&=x_2^2
  \quad\Longrightarrow\quad x_1=x_2=x_3,
  \qquad\text{fals.}
\end{aligned}
\]
; fără preambul sau pagină nouă.
% NOTE EDITORIALE:
% - La exercițiul 7, manuscrisul pare să scrie "nu e inj.".
%   Incluziunea canonică Z -> Q ESTE injectivă, însă nu este surjectivă.
%   Păstrăm lectura aparentă, marcată "ca în manuscris", fără corectare
%   tacită.
% - Monoizii din exercițiul 7 sunt priviți cu înmulțirea.
%   Diagrama arată definiția epimorfismului prin egalitatea a două
%   morfisme din Q care coincid pe Z.
% - În calculul pentru v, manuscrisul folosește o linie lungă în locul
%   repetării pașilor; păstrăm această prescurtare.
% - În exercițiul 8 se presupune un izomorfism ipotetic: funcția este
%   bijectivă și respectă produsul, iar la final se folosește injectivitatea.
% - Notația de lângă acolada pentru f(a)f(b)=4 are un al doilea rând
%   greu de citit (f(b)=...?); f(a)=1 => a=1 este clar.
% - "Numere prime" se referă la numere prime în N^*.
% - Textul șters și corecturile barate nu sunt reconstruite.

\subsection*{Seminar 4 - Pagina 3}

\noindent\underline{Exercițiul 7.}
Incluziunea
\[
 i:(\mathbb Z,\cdot)\longrightarrow(\mathbb Q,\cdot),
 \qquad i(x)=x,
\]
este epimorfism de monoizi care \emph{nu e injectiv}
(așa cum apare în manuscris).

\smallskip
\noindent\underline{Epimorfism.}
Pentru orice monoid $M$ și orice morfisme de monoizi $u,v$
astfel încât $u\circ i=v\circ i$, avem $u=v$.
\[
\begin{tikzcd}[column sep=large]
 \mathbb Z \arrow[r,"i"] &
 \mathbb Q \arrow[r,shift left=.7ex,"u"]
            \arrow[r,shift right=.7ex,swap,"v"] & M
\end{tikzcd}
\]

\noindent Fie $(M,\cdot)$ monoid și $u,v:\mathbb Q\to M$
morfisme de monoizi cu $u\circ i=v\circ i$. Atunci
\[
 u(m)=v(m)\qquad(\forall m\in\mathbb Z).
\]
Vrem să arătăm că $u(r)=v(r)$ pentru orice $r\in\mathbb Q$.
Fie $m\in\mathbb Z$, $m\ne0$. Avem
\[
\begin{aligned}
 e=u(1)
  &=u\!\left(\frac{m}{m}\right)
   =u\!\left(m\cdot\frac1m\right)
   =u(m)u\!\left(\frac1m\right)
   =u\!\left(\frac1m\right)u(m),\\
 e=v(1)
  &=\cdots
   =v\!\left(\frac1m\right)v(m).
\end{aligned}
\]
Deci $u(1/m)$ este inversul lui $u(m)$, iar $v(1/m)$ este
inversul lui $v(m)$. Cum $u(m)=v(m)$, rezultă
\[
 u\!\left(\frac1m\right)=v\!\left(\frac1m\right).
\]
Pentru $r=n/m\in\mathbb Q$, cu $m\ne0$,
\[
 u(r)=u(n)u\!\left(\frac1m\right)
     =v(n)v\!\left(\frac1m\right)=v(r).
\]
Prin urmare $u=v$.

\smallskip
\noindent [Orice epimorfism de grupuri este surjectiv.]

\medskip
\noindent\underline{Exercițiul 8.}
Monoizii
\[
 (\mathbb N^*,\cdot),\qquad
 (3\mathbb N+1,\cdot),\qquad
 3\mathbb N+1=\{3k+1\mid k\in\mathbb N\},
\]
nu sunt izomorfi.

\noindent\underline{Demonstrație.}
Presupunem că există un izomorfism
\[
 f:\mathbb N^*\longrightarrow 3\mathbb N+1,
 \qquad f(xy)=f(x)f(y)
 \quad(\forall x,y\in\mathbb N^*).
\]
Din surjectivitatea lui $f$ rezultă că există $x_1,x_2,x_3\in\mathbb N^*$
astfel încât
\[
 f(x_1)=4,\qquad f(x_2)=10,\qquad f(x_3)=25.
\]
Arătăm că $x_1,x_2,x_3$ sunt numere prime.
Dacă $x_1=ab$, cu $a,b>1$, atunci
\[
 f(x_1)=f(a)f(b)=4
 \quad\Longrightarrow\quad f(a)=1
 \quad\Longrightarrow\quad a=1,
 \qquad\text{fals.}
\]
Analog, $x_2$ și $x_3$ sunt numere prime.
Dar $4\cdot25=10^2$, deci
\[
\begin{aligned}
 f(x_1)f(x_3)&=f(x_2)^2,\\
 f(x_1x_3)&=f(x_2^2),\\
 x_1x_3&=x_2^2
  \quad\Longrightarrow\quad x_1=x_2=x_3,
  \qquad\text{fals.}
\end{aligned}
\]
